La empresa \textbf{MueblesPro} produce y comercializa tres tipos de productos: \textbf{mesas (M)}, \textbf{sillas (S)} y \textbf{estanterías (E)}, que son transportadas en lotes, y estos pueden estar completos o fraccionados (se puede enviar un lote completo, medio lote, un cuarto de lote, etc.). Por cada lote de estos productos, la empresa obtiene unos ingresos netos de \textbf{20, 15 y 25 unidades monetarias}, respectivamente.

La fábrica dispone de una planta con \textbf{80 horas de producción semanales}. Producir un lote de mesas lleva \textbf{15 horas}, un lote de sillas lleva \textbf{10 horas} y montar un lote de estanterías lleva \textbf{20 horas}.

Además, por contrato con sus distribuidores, la empresa debe suministrar al menos el contenido equivalente a \textbf{6 lotes completos en total}, entre mesas, sillas y estanterías.

El modelo de programación lineal que permite determinar el plan de producción óptimo es el siguiente:

\begin{equation}
\begin{split}
\text{max. } z = 20x_1 + 15x_2 + 25x_3 \\
\text{s.a.:} \\
15x_1 + 10x_2 + 20x_3 \leq 80 \quad \text{(tiempo de montaje disponible)} \\
x_1 + x_2 + x_3 \geq 6 \quad \text{(compromiso de producción total)} \\
x_1, x_2, x_3 \geq 0 \quad \text{(No negatividad)}
\end{split}
\end{equation}

Se pide

\begin{enumerate}

    \item \textbf{Obtener el plan de programación de la producción óptimo e indicar los niveles de producción, el beneficio y el cumplimiento de los requisitos descritos. Se puede hacer uso de la información de que, en la solución óptima, se cumple que $x_1=x_3=x_4=h_1=0$.} 

    El problema se puede formular en formato estándar:

    \begin{equation}
\begin{split}
\text{max. } z = 20x_1 + 15x_2 + 25x_3 \\
\text{s.a.:} \\
15x_1 + 10x_2 + 20x_3 + h_1 = 80 \quad \text{(tiempo de montaje disponible)} \\
x_1 + x_2 + x_3 - h_2 = 6 \quad \text{(compromiso de producción total)} \\
x_1, x_2, x_3, h_1, h_2 \geq 0 \quad \text{(No negatividad)}
\end{split}
\end{equation}

    Como el problema tiene 5 variables y 2 restricciones, existen tres variables no básicas y dos básicas. Dado que se sabe que $x_1=x_3=h_1=0$ en la solución óptima, se pueden seleccionar las variables $x_2$ y $h_2$ como variables básicas.

    En lugar de realizar todas las iteraciones con el método de las dos fases o con el método de la M grande, se puede obtener la tabla correspondiente a la solución óptima convirtiendo las columnas en $A$ de las variables $h_2$ y $x_2$ en columnas de la matriz identidad y en 0 las componentes correspondientes de esas variables en la función objetivo:


    \begin{center}
    \begin{tabular}{c|c|cccccc|c}
     & $z$ & $x_1$ & $x_2$ & $x_3$ & $h_1$ & $h_2$ & $a_2$ & \\   
     & 0 & 20 & 15 & 25 & 0 & 0 & 0 & \\ 
    \hline
    $h_1$ & 80  & 15 & 10 & 20 & 1 & 0 & 0 & $F_1$\\ 
    $a_2$ & 6  & 1 & 1 & 1 & 0 & -1 & 1& $F_2$\\ 
    \hline
    & & & & ... & & & &\\
    \hline
    Fase 2 & -120 & $-5/2$ & 0 & -5 & $-3/2$ & 0 & 0 & \\ 
    \hline
    $x_2$ & 8  & $3/2$ & 1 & 2 & $1/10$ & 0 & 0 & $F'_1 = F_1/10$\\ 
    $h_2$ & 2  & $1/2$ & 0 & 1 & $1/10$ & 1 & -1 & $F'_2 = -F_2 + F'_1$ \\ 
    \hline
    \end{tabular}
    \end{center}   

    
    Otra forma de obtener la tabla es conseguir la inversa de la base. Si elegimos como variables básicas $x_2$ y $h_2$, entonces la inversa de la base es

    \begin{equation}
    \begin{split}
    B^{-1} =\begin{pmatrix}
    1/10 & 0\\
    1/10 & -1\\
    \end{pmatrix}
    \end{split}
    \end{equation}

    Y con ello, recuperar los términos de la tabla $u_B$, $p_B$, $V^B$ y $z^B$

    \begin{equation}
\begin{split}
u^B=B^{-1}b = \begin{pmatrix}
1/10 & 0\\
1/10 & -1\\
\end{pmatrix}
\begin{pmatrix}
80\\
6\\
\end{pmatrix}
 = \begin{pmatrix}
8\\
2\\
\end{pmatrix}
\end{split}
\end{equation}

\begin{equation}
\begin{split}
p^B = B^{-1}A = \begin{pmatrix}
1/10 & 0\\
1/10 & -1\\
\end{pmatrix}
\begin{pmatrix}
15 & 10 & 20 & 1 & 0\\
1 & 1 & 1 & 0 & -1\\
\end{pmatrix}
 = \begin{pmatrix}
3/2 & 1 & 2 & 1/10 & 0\\
1/2 & 0 & 1 & 1/10 & 1\\
\end{pmatrix}
\end{split}
\end{equation}

\begin{equation}
\begin{split}
V^B = c-c^BB^{-1}A = \begin{pmatrix}
20 & 15 & 25 & 0 & 0\\
\end{pmatrix}
 - \begin{pmatrix}
15 & 0\\
\end{pmatrix}
\begin{pmatrix}
1/10 & 0\\
1/10 & -1\\
\end{pmatrix}
\begin{pmatrix}
15 & 10 & 20 & 1 & 0\\
1 & 1 & 1 & 0 & -1\\
\end{pmatrix}
 = \\
 \begin{pmatrix}
-5/2 & 0 & -5 & -3/2 & 0\\
\end{pmatrix}
\end{split}
\end{equation}

\begin{equation}
\begin{split}
z^B=c^Bu^B = \begin{pmatrix}
15 & 0\\
\end{pmatrix}
\begin{pmatrix}
8\\
2\\
\end{pmatrix}
 = 120\end{split}
\end{equation}

    
    Por último, se puede resolver el problema mediante el método de las dos fases.

    \begin{center}
    \begin{tabular}{c|c|cccccc|}
     & $z$ & $x_1$ & $x_2$ & $x_3$ & $h_1$ & $h_2$ & $a_2$\\ 
    Fase 1 & 6 & 1 & 1 & 1 & 0 & -1 & 0 \\ 
    Fase 2 & 0 & 20 & 15 & 25 & 0 & 0 & 0 \\ 
    \hline
    $h_1$ & 80  & 15 & 10 & 20 & 1 & 0 & 0\\ 
    $a_2$ & 6  & 1 & 1 & 1 & 0 & -1 & 1\\ 
    \hline
    Fase 1 & $2/3$ & 0 & $1/3$ & $-1/3$ & $-1/15$ & -1 & 0 \\ 
    Fase 2 & $-320/3$ & 0 & $5/3$ & $-5/3$ & $-4/3$ & 0 & 0 \\ 
    \hline
    $x_1$ & 16/3  & 1 & $2/3$ & $4/3$ & $1/15$ & 0 & 0\\ 
    $a_2$ & 2/3  & 0 & $1/3$ & $-1/3$ & $-1/15$ & -1 & 1\\ 
    \hline
    Fase 1 & 0 & 0 & 0 & 0 & 0 & 0 & -1 \\ 
    Fase 2 & -110 & 0 & 0 & 0 & -1 & 5 & -5 \\ 
    \hline
    $x_1$ & 4  & 1 & 0 & 2 & $1/5$ & 2 & -2\\ 
    $x_2$ & 2  & 0 & 1 & -1 & $-1/5$ & -3 & 3\\ 
    \hline
    Fase 2 & -120 & $-5/2$ & 0 & -5 & $-3/2$ & 0 & 0 \\ 
    \hline
    $h_2$ & 2  & $1/2$ & 0 & 1 & $1/10$ & 1 & -1\\ 
    $x_2$ & 8  & $3/2$ & 1 & 2 & $1/10$ & 0 & 0\\ 
    \hline
    \end{tabular}
    \end{center}

    La última tabla se podría escribir cambiando el orden para que su orden coincida con el orden que se ha seleccionado en las otras dos alternativas.

    \begin{center}
    \begin{tabular}{c|c|cccccc|}
     & $z$ & $x_1$ & $x_2$ & $x_3$ & $h_1$ & $h_2$ & $a_2$\\ 
      & -120 & $-5/2$ & 0 & -5 & $-3/2$ & 0 & 0 \\ 
    \hline
    $x_2$ & 8  & $3/2$ & 1 & 2 & $1/10$ & 0 & 0\\ 
    $h_2$ & 2  & $1/2$ & 0 & 1 & $1/10$ & 1 & -1\\ 
    \hline
    \end{tabular}
    \end{center}

    
En cualquiera de los caso, el plan de programación óptimo consiste en:
\begin{itemize}
    \item Producir 8 lotes de sillas.
    \item No producir ni mesas ni estanterías.
    \item Hacer uso al completo de las 80 horas de montaje.
    \item Cumplir el compromiso comercial superando en dos lotes el compromiso mínimo.
\end{itemize}

\item \textbf{Describir cuál sería el nuevo plan de producción si fuera necesario entregar al menos 9 lotes en lugar de 6.}

   
Este cambio daría lugar a un cambio en b: $b=(80 \;\; 9)^T$, lo cual daría lugar a un cambio en $u^B=B^{-1}b$ y en $z^B=c^Bu^B$.

\begin{equation}
    u^B=B^{-1}b =
    \begin{pmatrix}
    1/10 & 0\\
    1/10 & -1\\
    \end{pmatrix}
    \begin{pmatrix}
    80\\
    9\\
    \end{pmatrix} = 
    \begin{pmatrix}
    8\\
    -1\\
    \end{pmatrix} 
\end{equation}

Al sustituir los valores en la tabla de la solución óptima se observa que la solución deja de ser factible y que no existen tasas de sustitución de las variables no básicas con respecto a $h_2$ (negativa) y, por lo tanto, no existe solución factible. 

    \begin{center}
    \begin{tabular}{c|c|cccccc|}
     & $z$ & $x_1$ & $x_2$ & $x_3$ & $h_1$ & $h_2$ & $a_2$\\ 
      & $z^B$ & $-5/2$ & 0 & -5 & $-3/2$ & 0 & 0 \\ 
    \hline
    $x_2$ & 8  & $3/2$ & 1 & 2 & $1/10$ & 0 & 0\\ 
    $h_2$ & -1 & $1/2$ & 0 & 1 & $1/10$ & 1 & -1\\ 
    \hline
    \end{tabular}
    \end{center}
Es decir, no sería posible cumplir con el nuevo compromiso de entregas.

\item \textbf{Calcular el rango de valores de la contribución unitaria al beneficio de las sillas dentro del cual resulta interesante mantener la base de la solución óptima.}

Para responder a esta pregunta es necesario realizar el análisis de sensibilidad de $c_2$ que, en la situación original, es 15.

\begin{equation}
\begin{split}
V^B = c-c^BB^{-1}A = \begin{pmatrix}
20 & c_2 & 25 & 0 & 0\\
\end{pmatrix}
 - \begin{pmatrix}
c_2 & 0\\
\end{pmatrix}
\begin{pmatrix}
1/10 & 0\\
1/10 & -1\\
\end{pmatrix}
\begin{pmatrix}
15 & 10 & 20 & 1 & 0\\
1 & 1 & 1 & 0 & -1\\
\end{pmatrix}
 = \\
 \begin{pmatrix}
20 & c_2 & 25 & 0 & 0\\
\end{pmatrix}
 - \begin{pmatrix}
c_2 & 0\\
\end{pmatrix}
\begin{pmatrix}
3/2 & 1 & 2 & 1/10 & 0\\
1/2 & 0 & 1 & 1/10 & 1\\
\end{pmatrix}=\\
 \begin{pmatrix}
20 & c_2 & 25 & 0 & 0\\
\end{pmatrix}-
\begin{pmatrix}
3/2c_2 & c_2 & 2c_2 & 1/10c_2 & 0
\end{pmatrix}=\\
 \begin{pmatrix}
20-3/2c_2 & 0 & 25-2c_2 & -1/10c_2 & 0\\
\end{pmatrix}\leq 0 \implies \\
\left\{
\begin{matrix}
20-3/2c_2 \leq 0 \implies c_2 \geq 40/3 \\
25-2c_2 \leq 0 \implies c_2 \geq 25/2\\
-1/10c_2 \leq 0  \implies c_2 \geq 0 
\end{matrix}
\right\} \implies c_2 \in [40/3, \infty)
\end{split}
\end{equation}

\item \textbf{Explicar el significado del criterio del simplex de la variable $x_1$ interpretando explícitamente el significado de las tasas de sustitución de esta variable. La explicación debe referirse al problema de MueblesPro (y no solo a los aspectos matemáticos del modelo de programación lineal planteado).}


\begin{equation}
\begin{split}
V^B_{x_1} = c_{x_1}-c^BB^{-1}A_{x_1} = c_{x_1}-c^Bp^B_{x_1}  = 20
 - \begin{pmatrix}
15 & 0\\
\end{pmatrix}
\begin{pmatrix}
3/2 \\ 
1/2 
\end{pmatrix}
 = 
-5/2 
\end{split}
\end{equation}

Realizar una mesa ($x_1 = 1$) daría lugar a:
\begin{itemize}
    \item Un ingreso neto de 20 unidades por la venta de esa mesa ($c_{x_1}$)
    \item Una producción de 3/2 sillas menos ($p^B_{11}=3/2$).
    \item Una reducción de los ingresos por importe igual a 45/2: por cada unidad que se deja de producir y vender se dejan de ingresar 15 euros ($c_{x_2}$) y como se dejan de fabricar y vender 3/2 sillas, se dejan de ingresar $c_{x_2}p^B_{11}=15\times 3/2 = 45/2$.
    \item Una disminución de 1/2 de $h_2$ ($p_{21}=1/2$), lo cual representa una disminución de 1/2 de la medida en la que el compromiso de entregas se está cumpliendo por encima del mínimo acordado; pero esto no tiene impacto en la función objetivo ($c_{h_2}$=0).
\end{itemize}

Como resultado de todo lo anterior, fabricar una mesa daría lugar a un ingreso de 20 por su venta y una reducción de los ingresos de 45/2 por la reducción de ventas de sillas, lo que globlamente da lugar a una reducción de 5/2 del beneficio total ($V^B_{x_2}=-5/2$).
        
\end{enumerate}

